\chapter{Conclusion}
\label{sec:conclusion}
% Answer each research question and address how the limitations of the study qualify the conclusion.
% Template guidance: make the relation between the research gap and the contribution
% clear, and be honest about how limitations qualify each answer.

%%%%%%%%%%%%%%%%%%%%%%%%%%%%%%%%%%%%%%%%%%%%%%%%%%%%%%%%%%%%%%%%%%%%%%%%%%%%%%%%
\section{Answers to the Research Questions}
\label{sec:conclusion:answers}
% One direct answer each -- a sentence or two, not a summary of the Results chapter.
% Where a limitation qualifies the answer, say so in the same breath.

\subsection{RQ1: Task Feasibility \& Baseline}
\label{sec:conclusion:rq1}

\subsection{RQ2: Reduction Efficiency}
\label{sec:conclusion:rq2}

\subsection{RQ3: Predictive Retention \& Trade-offs}
\label{sec:conclusion:rq3}

\subsection{RQ4: Value of Domain-Informed Adaptation}
\label{sec:conclusion:rq4}

\subsection{RQ5: Reduced to Unreduced? Generalization}
\label{sec:conclusion:rq5}


%%%%%%%%%%%%%%%%%%%%%%%%%%%%%%%%%%%%%%%%%%%%%%%%%%%%%%%%%%%%%%%%%%%%%%%%%%%%%%%%
\section{Contributions Revisited}
\label{sec:conclusion:contributions}
% Close the loop with \ref{sec:intro:contributions} and the gap stated in
% \ref{sec:relwork:gap}: what was claimed, and what the evidence actually supports.


%%%%%%%%%%%%%%%%%%%%%%%%%%%%%%%%%%%%%%%%%%%%%%%%%%%%%%%%%%%%%%%%%%%%%%%%%%%%%%%%
\section{Future Work}
\label{sec:conclusion:futurework}
% Picks up the threads hinted at in the Discussion.

\subsection{Extending Beyond a Single Synthesis Algorithm}
\label{sec:conclusion:futurework:algorithms}
% The Deepsyn/Syn4/C2RS graphs already exist on disk; algorithm-ranking was the
% original motivation and is the most natural next step.

\subsection{Stronger Domain-Aware Reduction}
\label{sec:conclusion:futurework:domain}
% Especially if RQ4 came out null: what a genuinely strong AIG-aware method would
% need to do differently.

\subsection{Scaling to Industrial AIGs}
\label{sec:conclusion:futurework:scale}
% Removing the bounded-scale constraint, which is the limitation closest to the
% original motivation.

\subsection{From Prediction to Script Generation}
\label{sec:conclusion:futurework:generation}
% Closing the loop the motivation opens with: using predicted optimizability to
% drive data-driven script construction.
